\scanpage{12}
\mspage{11}

% vertical rule to the left of the condition
\uline{Fig}$_{\unsure{2|n}}$ \inserted{Soit $\deleted{F}\unsure{X|F|\mathbf{X}}\in\mathcal{F}_{el}$,} Considérons\vadjust{\vbox to 0pt{\hbox{\kern-0.7em\vrule width 0.35pt height 0pt depth 3\baselineskip}\vss}} l'ens ordonné des $\deleted{F}X'\leq\deleted{F}X$ est de \uline{dim comb. finie}, i.e. il existe un entier $n$, tel que tte partie finie $\deleted{\unsure{F}}\Phi$ \deleted{\illegible{1}} majorée par $X$, \inserted{et tot. ordonnée} soit de cardinal $\leq n$.

La \uline{\unsure{dimension|dimention}} \uline{combinatoire} d'un $X\in\mathcal{F}_{el}$ est le Sup des cardinaux des parties finies $\Phi$ précédentes, diminués de 1. Donc la dim combinatoire est $\geq 0$, et elle est nulle ssi $X$ est un \uline{élémt} \uline{minimal} de $\mathcal{F}_{el}$.

Dans les exemples \unsure{types|type}, \deleted{\unsure{F|les}} les figures élémentaires minimales correspondent aux parties fermées de $X$ qui sont des \uline{variétés} \uline{connexes}.

Si $\Phi\subset\mathcal{F}_{\text{él}}$ est une figure, \deleted{sa} \uline{dim combinatoire} est sa dim comb en tant qu'ens ordonné, c'est \unsure{ou|donc} le sup des $\mathrm{dim\,comb}\,X$, pour $X\in\Phi$. Elle est nulle ssi \deleted{$\Phi$ est formé de va} \inserted{$\Phi\neq\varnothing$ et} la relation
